\chapter{Sviluppo agile}
\label{cap:agile}

\begin{inbreve}
I prodotti devono arrivare presto sul mercato (capitolo~\ref{cap:prodotti}). Per questo si usano i metodi agili,
che consegnano il software a incrementi e riducono il lavoro di pianificazione e documentazione. Il capitolo
confronta lo sviluppo agile con quello guidato dal piano e riporta il manifesto e i dodici principi. Poi
descrive lo sviluppo incrementale e due metodi agili: Extreme Programming e, più a fondo, Scrum.
\end{inbreve}

\section{Sviluppo guidato dal piano e sviluppo agile}
\label{sec:agile-confronto}

Nel Novecento lo sviluppo era \termine{plan-driven}. Si basava su:
\begin{itemize}
  \item pianificazione dettagliata del progetto;
  \item specifica dei requisiti;
  \item metodi di analisi e di design;
  \item documentazione completa del sistema;
  \item controllo formale della qualità.
\end{itemize}

Questo approccio comporta un \ingl{overhead} importante di pianificazione, progettazione e documentazione. Per i
prodotti software di dimensioni piccole e medie l'overhead finisce per dominare il processo di sviluppo.

\begin{definizione}[Agile software engineering]
L'\termine{Agile SE} concentra l'attenzione sul software, per:
\begin{itemize}
  \item consegnare le funzionalità rapidamente;
  \item rispondere rapidamente ai cambiamenti;
  \item minimizzare l'overhead di sviluppo.
\end{itemize}
\end{definizione}

\section{Il manifesto agile}
\label{sec:agile-manifesto}

Il \termine{Manifesto for Agile Software Development} (2001, \texttt{agilemanifesto.org}) dichiara quattro valori.
Ogni valore ha due parti: c'è valore in entrambe, ma si preferisce quella a sinistra.

\begin{center}
\begin{tabular}{@{}rcl@{}}
Individui e interazioni & \emph{più di} & processi e strumenti \\
Software funzionante & \emph{più di} & documentazione completa \\
Collaborazione con il cliente & \emph{più di} & negoziazione del contratto \\
Risposta al cambiamento & \emph{più di} & aderenza a un piano \\
\end{tabular}
\end{center}

Il manifesto è sviluppato in dodici principi.

\begin{enumerate}
  \item La priorità massima è soddisfare il cliente con consegne di software di valore, rapide e continue.
  \item I cambiamenti dei requisiti sono benvenuti, anche a sviluppo avanzato. I processi agili li sfruttano a
    vantaggio competitivo del cliente.
  \item Si consegna software funzionante frequentemente, da un paio di settimane a un paio di mesi, preferendo
    la scala di tempo più breve.
  \item Persone di business e sviluppatori lavorano insieme ogni giorno per tutto il progetto.
  \item I progetti si costruiscono attorno a persone motivate. Si dà loro l'ambiente e il supporto di cui hanno bisogno
    e si ha fiducia che portino a termine il lavoro.
  \item Il modo più efficiente ed efficace per comunicare con e dentro un team è la conversazione faccia a faccia.
  \item Il software funzionante è la misura principale del progresso.
  \item I processi agili promuovono uno sviluppo sostenibile: sponsor, sviluppatori e utenti devono poter mantenere
    un ritmo costante indefinitamente.
  \item L'attenzione continua all'eccellenza tecnica e a un buon design aumenta l'agilità.
  \item La semplicità, cioè l'arte di massimizzare il lavoro \emph{non} fatto, è essenziale.
  \item Le migliori architetture, requisiti e design emergono da team che si auto-organizzano.
  \item A intervalli regolari il team riflette su come diventare più efficace e di conseguenza adatta il proprio
    comportamento.
\end{enumerate}

\section{Sviluppo incrementale}
\label{sec:agile-incrementale}

Tutti i metodi agili si basano sullo \termine{sviluppo e consegna incrementali}. Il prodotto è visto come un
insieme di \termine{feature} (\rif{sec:feature-intro}) e viene costruito aggiungendone alcune alla volta
(\figrif{fig:agile-incrementale}).

\begin{figure}[ht]
\centering
\begin{tikzpicture}[blocco/.append style={minimum width=3.1cm, minimum height=1cm}]
\node[blocco] (a) at (0,0) {Choose features to be\\included in increment};
\node[blocco] (b) at (5.2,0) {Refine feature\\descriptions};
\node[blocco] (c) at (5.2,-2.1) {Implement and\\test feature};
\node[blocco] (d) at (0,-2.1) {Integrate feature\\into system};
\node[blocco] (e) at (-5.2,-2.1) {Deliver system\\increment};
\node[font=\small, align=center] (l) at (0,1.6) {Product feature list};
\node[font=\small, align=center] (r) at (-5.2,-3.9) {Se tutte le feature sono\\complete: system release};
\draw[freccia] (l) -- (a);
\draw[freccia] (a) -- (b);
\draw[freccia] (b) -- (c);
\draw[freccia] (c) -- (d);
\draw[freccia] (d) -- (e);
\draw[freccia] (e.north) |- (a.west);
\draw[freccia] (e) -- (r);
\end{tikzpicture}
\caption{Ciclo dello sviluppo incrementale. Dall'elenco delle feature si scelgono quelle dell'incremento, se ne
raffinano le descrizioni, si implementano e si collaudano, si integrano nel sistema e si consegna l'incremento. Si
riparte finché tutte le feature sono complete: allora si consegna la release del sistema.}
\label{fig:agile-incrementale}
\end{figure}

\section{Extreme Programming}
\label{sec:agile-xp}

\termine{Extreme Programming} (XP) è un metodo agile che riunisce dieci pratiche:
\ingl{incremental planning}, \ingl{small releases}, \ingl{simple design}, \ingl{test-first development},
\ingl{refactoring}, \ingl{pair programming}, \ingl{collective ownership}, \ingl{continuous integration},
\ingl{sustainable pace}, \ingl{on-site customer}. Cinque di esse sono descritte nella \tabrif{tab:agile-xp}.
Le altre cinque sono spiegate subito dopo la tabella.

\begin{table}[ht]
\centering
\small
\caption{Cinque pratiche di Extreme Programming.}
\label{tab:agile-xp}
\begin{tabularx}{\linewidth}{>{\raggedright\arraybackslash}p{3.4cm}X}
\toprule
Pratica & Descrizione \\
\midrule
Incremental planning / user stories & Non c'è un ``piano generale''. Ciò che va implementato (i requisiti) in
  ogni incremento si stabilisce in discussioni con un rappresentante del cliente. I requisiti sono scritti come
  user stories (\rif{sec:feature-storie}). Le storie incluse in una release dipendono dal tempo disponibile e
  dalla loro priorità relativa. \\
Small releases & Si sviluppa per primo il minimo insieme utile di funzionalità che dà valore di business. Le release
  sono frequenti e aggiungono funzionalità alla release precedente. \\
Test-driven development & Invece di scrivere il codice e poi i test, gli sviluppatori scrivono prima i test.
  Questo chiarisce che cosa deve fare il codice e garantisce che esista sempre una versione ``testata''. Un
  framework di unit test automatici esegue i test a ogni modifica. Il codice nuovo non deve ``rompere'' quello già
  implementato. \\
Continuous integration & Appena un task è completo, viene integrato nell'intero sistema e si crea una nuova versione.
  Gli unit test di tutti gli sviluppatori sono eseguiti automaticamente e devono avere successo prima che la nuova
  versione sia accettata. \\
Refactoring & Significa migliorare struttura, leggibilità, efficienza e sicurezza di un programma. Tutti gli
  sviluppatori devono fare refactoring del codice appena trovano possibili miglioramenti. Questo mantiene il
  codice semplice e manutenibile. \\
\bottomrule
\end{tabularx}
\end{table}

\subsection{Le altre cinque pratiche\aggiunta}

\begin{description}[leftmargin=0pt, style=unboxed]
  \item[Simple design.] Si progetta la soluzione più semplice che soddisfa i requisiti di \emph{adesso}, senza
    prevedere esigenze future. Se i requisiti cambiano, il refactoring adatta il codice. È l'applicazione del principio
    della semplicità (\rif{sec:agile-manifesto}, principio 10).
  \item[Pair programming.] Due sviluppatori lavorano insieme alla stessa postazione. Uno scrive il codice, l'altro lo
    legge e lo controlla subito. Il codice è rivisto mentre nasce e la conoscenza si diffonde nel team.
  \item[Collective ownership.] Il codice non appartiene a una persona. Ogni sviluppatore può modificare qualsiasi
    parte del sistema. Gli unit test proteggono da modifiche che rompono qualcosa.
  \item[Sustainable pace.] Si evita il lavoro straordinario continuo: un ritmo che non si può sostenere abbassa la
    qualità. Corrisponde al principio 8 del manifesto.
  \item[On-site customer.] Un rappresentante del cliente è presente con il team per tutto lo sviluppo. Risponde subito
    ai dubbi sui requisiti e decide le priorità. Corrisponde al principio 4 del manifesto.
\end{description}

\section{Scrum}
\label{sec:agile-scrum}

\begin{definizione}[Scrum]
\termine{Scrum} è un \ingl{framework} per organizzare un progetto agile. Prevede persone designate che fanno da
interfaccia tra il team di sviluppo e il resto dell'organizzazione.
\end{definizione}

\begin{table}[ht]
\centering
\small
\caption{Terminologia di Scrum.}
\label{tab:agile-scrum}
\begin{tabularx}{\linewidth}{>{\raggedright\arraybackslash}p{4.2cm}X}
\toprule
Termine & Significato \\
\midrule
Development team & Piccolo team che si auto-organizza (4--8 persone), responsabile dello sviluppo del prodotto. \\
Product Owner & Membro del team responsabile di identificare feature e attributi del prodotto, di rivedere il lavoro
  svolto e di aiutare a collaudarlo. \\
ScrumMaster & ``Coach'' del team: lo guida a usare Scrum in modo efficace. \\
Product backlog & Elenco delle voci (bug, feature, miglioramenti) che il team non ha ancora completato. \\
Sprint & Breve periodo (2--4 settimane) in cui si sviluppa un incremento del prodotto. \\
Potentially shippable product increment & Risultato di uno sprint, di qualità sufficiente per essere messo in uso dai
  clienti. \\
Scrum & Riunione quotidiana del team in cui si rivede l'avanzamento e si concorda il lavoro della giornata. \\
\bottomrule
\end{tabularx}
\end{table}

\subsection*{Il ciclo di Scrum}

Il prodotto è sviluppato in una serie di sprint (\figrif{fig:agile-scrum}):
\begin{itemize}
  \item le voci a priorità più alta del product backlog sono scelte per lo sprint successivo;
  \item la pianificazione dello sprint (\ingl{plan sprint}) crea lo \ingl{sprint backlog};
  \item nello sprint il team implementa quante più voci dello sprint backlog può nel tempo dato. Le voci incomplete
    tornano nel product backlog;
  \item la riunione di revisione serve a imparare gli uni dagli altri come evitare i problemi e migliorare la
    produttività negli sprint successivi.
\end{itemize}
L'ultimo punto richiama il \emph{Toyota Kata}: favorire il miglioramento continuo dando ai dipendenti il potere di
individuare i problemi e di trovare insieme le soluzioni.

\begin{figure}[ht]
\centering
\begin{tikzpicture}[blocco/.append style={minimum width=3cm, minimum height=0.9cm}]
\node[blocco, fill=grigiochiaro] (pb) at (0,0) {Product backlog};
\node[blocco] (sel) at (4.7,0) {Select items\\to implement};
\node[blocco] (pl) at (9.4,0) {Plan sprint};
\node[blocco, fill=grigiochiaro] (sb) at (9.4,-2.2) {Sprint backlog};
\node[blocco] (sp) at (4.7,-2.2) {Sprint\\(develop software,\\test software, scrum)};
\node[blocco] (inc) at (0,-2.2) {Shippable\\product increment};
\node[blocco] (rev) at (0,-4.4) {Review sprint};
\draw[freccia] (pb) -- (sel);
\draw[freccia] (sel) -- (pl);
\draw[freccia] (pl) -- (sb);
\draw[freccia] (sb) -- (sp);
\draw[freccia] (sp) -- (inc);
\draw[freccia] (inc) -- (rev);
\draw[freccia] (rev.west) -- (-2.6,-4.4) |- (pb.west);
\node[font=\footnotesize, align=right, anchor=east] at (-2.75,-2.2) {review\\product\\backlog};
\end{tikzpicture}
\caption{Ciclo di Scrum. Le voci scelte dal product backlog diventano lo sprint backlog. Lo sprint produce un
incremento potenzialmente consegnabile. La revisione dello sprint alimenta il product backlog per lo sprint
successivo.}
\label{fig:agile-scrum}
\end{figure}

\section{Il product backlog}
\label{sec:agile-backlog}

Il \termine{product backlog} è l'elenco di ciò che resta da fare per completare lo sviluppo del prodotto. Le sue
voci si chiamano \termine{product backlog item} (PBI). Il backlog deve essere sempre \emph{ordinato per priorità}.

\begin{esempio}[Un product backlog]
\begin{enumerate}
  \item As a teacher, I want to be able to configure the group of tools that are available to individual classes.
    (\emph{feature})
  \item As a parent, I want to be able to view my children's work and the assessments made by their teachers.
    (\emph{feature})
  \item As a teacher of young children, I want a pictorial interface for children with limited reading ability.
    (\emph{user request})
  \item Establish criteria for the assessment of open source software that might be used as a basis for parts of
    this system. (\emph{development activity})
  \item Refactor user interface code to improve understandability and performance. (\emph{engineering improvement})
  \item Implement encryption for all personal user data. (\emph{engineering improvement})
\end{enumerate}
Le voci 1--3 sono \emph{user stories} (\rif{sec:feature-storie}), la 4 è un'attività del team, le 5 e 6 sono
miglioramenti del software.
\end{esempio}

Ogni PBI si trova in uno di tre stati.

\begin{description}[leftmargin=0pt, style=unboxed]
  \item[Ready for consideration.] Idee e descrizioni di feature di alto livello, da considerare per l'inclusione nel
    prodotto. Sono provvisorie: possono cambiare radicalmente o non essere incluse nel prodotto finale.
  \item[Ready for refinement.] Il team ha concordato che la voce è importante e va implementata nello sviluppo
    corrente. La definizione di ciò che serve è ragionevolmente chiara, ma serve lavoro per capirla e raffinarla.
  \item[Ready for implementation.] La voce ha abbastanza dettaglio perché il team ne stimi lo sforzo e la
    implementi. Le dipendenze da altre voci sono state identificate.
\end{description}

Il backlog è mantenuto con quattro attività.

\begin{description}[leftmargin=0pt, style=unboxed]
  \item[Refinement.] Le voci esistenti sono analizzate e raffinate in voci più dettagliate. Questo può far nascere
    nuove voci.
  \item[Estimation.] Il team stima il lavoro necessario a implementare un PBI e lo aggiunge a ogni voce analizzata.
    Lo sforzo si esprime in \emph{person-hour} o \emph{person-day}, cioè le ore o i giorni che servirebbero a una
    persona. Per esempio, 10 person-day con 2 persone richiedono 5 giorni.
  \item[Creation.] Si aggiungono voci nuove: feature suggerite dal product manager, modifiche richieste alle feature,
    miglioramenti tecnici o attività di processo (per esempio valutare strumenti di sviluppo).
  \item[Prioritization.] Le voci sono riordinate per tenere conto di nuove informazioni e di circostanze cambiate.
\end{description}

\section{Sprint con tempo fisso}
\label{sec:agile-sprint}

Il prodotto è sviluppato in sprint, ciascuno dei quali consegna un incremento. Gli sprint durano poco (1--4
settimane; la terminologia di Scrum, \tabrif{tab:agile-scrum}, indica 2--4) e hanno una data di inizio e una di fine definite. Sono \termine{timeboxed}: lo sviluppo si ferma alla
fine dello sprint, che il lavoro sia completo o no. Durante lo sprint il team lavora sulle voci del product backlog.

\subsection*{Benefici del timeboxing}

\begin{itemize}
  \item \emph{Progresso dimostrabile}: a ogni sprint c'è un risultato tangibile (di solito un dimostratore
    software) che si può consegnare.
  \item \emph{Problem discovery}: se si scoprono errori o omissioni, il lavoro da rifare è limitato alla durata di
    uno sprint.
  \item \emph{Work planning}: il team capisce quanto lavoro riesce a fare in un periodo fisso.
\end{itemize}

\subsection*{Le attività di uno sprint}

\begin{description}[leftmargin=0pt, style=unboxed]
  \item[Sprint planning.] Si scelgono le voci da completare nello sprint e, se serve, si raffinano per creare lo
    sprint backlog. Non dovrebbe durare più di un giorno, all'inizio dello sprint. Comprende tre passi:
    \begin{itemize}
      \item fissare un \emph{sprint goal} concordato (può riguardare funzionalità, supporto, prestazioni o
        affidabilità);
      \item decidere l'elenco dei PBI da implementare. Gli input chiave sono le stime di sforzo associate ai PBI
        e la \termine{velocity} del team, cioè la quantità di lavoro che il team riesce a completare in uno sprint;
      \item creare lo \emph{sprint backlog}, una versione più dettagliata del product backlog che registra il lavoro
        da fare nello sprint.
    \end{itemize}
  \item[Sprint execution.] Il team implementa le voci dello sprint backlog scelte. Se non riesce a completarle
    tutte, lo sprint \emph{non} viene allungato: le voci non finite tornano nel product backlog e sono messe in coda per
    uno sprint futuro.
  \item[Scrum quotidiano.] Durante lo sprint il team ha riunioni giornaliere per coordinare il lavoro. È breve
    (\ingl{stand-up}) e mirata, di solito all'inizio della giornata. Ognuno condivide informazioni, descrive i
    progressi dal giorno prima, i problemi sorti e i piani per il giorno. Così tutti sanno cosa succede e, se nascono
    problemi, possono ripianificare il lavoro a breve termine. Si rivede lo sprint backlog: le voci completate
    sono tolte e, se emergono nuove informazioni, se ne aggiungono altre. Poi il team decide chi lavora su quali
    voci quel giorno.
  \item[Sprint reviewing.] Il lavoro svolto è rivisto dal team e, eventualmente, da stakeholder esterni. Il team
    riflette anche su cosa è andato bene e cosa è andato male, per migliorare il proprio modo di lavorare.
\end{description}

Una voce si considera completa quando soddisfa una \emph{code completeness checklist}. Un esempio:

\begin{itemize}
  \item \emph{Reviewed}: il codice è stato rivisto da un altro membro del team, che ha controllato il rispetto degli
    standard di codifica, la comprensibilità, i commenti adeguati e l'eventuale refactoring.
  \item \emph{Unit tested}: tutti gli unit test sono stati eseguiti automaticamente e con successo.
  \item \emph{Integrated}: il codice è integrato nel codebase del progetto e non sono stati segnalati errori di
    integrazione.
  \item \emph{Integration tested}: tutti gli integration test sono stati eseguiti automaticamente e con successo.
  \item \emph{Accepted}: se opportuno, sono stati eseguiti gli acceptance test e il product owner o il team di
    sviluppo hanno confermato che la voce è completa.
\end{itemize}

Scrum non prescrive le attività tecniche agili da usare. Due pratiche però vanno \emph{sempre} usate in uno sprint:
la \termine{test automation} e la \termine{continuous integration}.

\subsection{Test automation e continuous integration\aggiunta}

La \emph{test automation} consiste nello scrivere test che il computer esegue da solo, senza intervento manuale,
per esempio dopo ogni modifica del codice. Così un errore nuovo viene scoperto subito, mentre il lavoro è ancora
fresco. La \emph{continuous integration} richiede di integrare spesso il lavoro di ciascuno nel codebase comune, anche
più volte al giorno. A ogni integrazione partono i test automatici. Le due pratiche si completano: senza test
automatici l'integrazione frequente sarebbe rischiosa, perché non si saprebbe se la modifica ha rotto qualcosa.

\begin{esempio}[Pianificare uno sprint\aggiunta]
Un team di 5 persone ha uno sprint di 2 settimane, cioè 10 giorni di lavoro ciascuno: la capacità è
$5 \times 10 = 50$ person-day. Il backlog, già ordinato per priorità, contiene voci stimate 8, 15, 12, 10 e 9
person-day. Si prendono le voci in ordine finché la somma sta nella capacità: $8+15 = 23$, $23+12 = 35$,
$35+10 = 45$. La voce successiva porterebbe la somma a $45+9 = 54 > 50$: resta nel product backlog per lo
sprint seguente. In pratica la capacità si stima dalla \emph{velocity}, cioè dal lavoro realmente completato negli
sprint passati, non dal calendario.
\end{esempio}

\section{Team che si auto-organizzano}
\label{sec:agile-team}

\subsection*{Composizione}

Un team Scrum affronta compiti diversi (per esempio rete, esperienza utente, progettazione di database) e quindi
richiede persone con competenze ed esperienza diverse. Un team di 5--8 persone è abbastanza grande per essere vario
e abbastanza piccolo per comunicare in modo informale ed efficace e per concordare le priorità. Poiché è il team, e
non il singolo, a rispondere del lavoro, è più facile gestire chi esce o entra. Con una buona comunicazione ognuno
impara qualcosa sulle aree degli altri.

\subsection*{Coordinamento}

Gli ideatori di Scrum supponevano team nella stessa sede. Lo scrum quotidiano si fonda però su due assunzioni che
non sempre valgono:
\begin{itemize}
  \item il team è fatto di lavoratori a tempo pieno che condividono lo spazio di lavoro. In realtà i membri possono
    essere part-time e lavorare in luoghi diversi;
  \item tutti possono partecipare a una riunione mattutina. In realtà alcuni hanno orari flessibili o lavorano a
    più progetti insieme.
\end{itemize}

\subsection*{Interazioni esterne}

Le \termine{interazioni esterne} sono quelle dei membri del team con persone fuori dal team. Solo lo ScrumMaster e il
Product Owner dovrebbero occuparsene: gli sviluppatori si concentrano sullo sviluppo.

\subsection*{Project management}

In ogni azienda i team devono riferire l'avanzamento alla direzione. Gli ideatori di Scrum non avevano previsto che lo
ScrumMaster avesse anche responsabilità di project management. In molte aziende però succede, perché lo ScrumMaster
conosce il lavoro in corso ed è nella posizione migliore per fornire informazioni accurate, piani di progetto e
avanzamento.

\begin{puntichiave}
\begin{itemize}
  \item Plan-driven: overhead dominante per prodotti piccoli e medi. Agile: consegnare presto, rispondere ai
    cambiamenti, ridurre l'overhead.
  \item Il manifesto dà quattro valori e dodici principi. Il software funzionante è la misura principale del progresso.
  \item Sviluppo incrementale: si scelgono feature, si raffinano, si implementano, si integrano e si consegna un
    incremento.
  \item XP: pratiche come test-driven development, continuous integration e refactoring.
  \item Scrum: team di 4--8 persone, Product Owner, ScrumMaster, product backlog, sprint di 2--4 settimane,
    scrum quotidiano.
  \item Lo sprint è timeboxed: non si allunga, le voci non finite tornano nel backlog. Test automation e continuous
    integration sono sempre usate.
\end{itemize}
\end{puntichiave}

\begin{domande}
\begin{enumerate}
  \item Perché l'overhead del plan-driven development pesa sui prodotti di dimensioni piccole e medie?
  \item Quali sono i quattro valori del manifesto agile?
  \item Che cosa significa che gli sprint sono timeboxed e quali sono i benefici?
  \item Quali ruoli ha Scrum e quali sono le responsabilità di ciascuno?
  \item In quali stati può trovarsi un PBI e quali attività mantengono il backlog?
  \item Quali assunzioni sul team fa lo scrum quotidiano e quando non valgono?
  \item Quali pratiche tecniche vanno sempre usate in uno sprint?
\end{enumerate}
\end{domande}
